\section{Conclusion}
\label{sec:ch6_conclusion}

This dissertation has advanced the state of the art in edge-cloud orchestration by demonstrating that computational efficiency, cross-application generalization, and validated autonomy can be achieved simultaneously.

\ac{AERO} proved that domain-specific design enables edge-deployable prediction with only 599 parameters, representing a 99.98\% reduction from transformer baselines, while achieving sub-millisecond inference and 99\% fewer \ac{SLA} violations. \ac{OmniFORE} demonstrated that temporal clustering and multi-scale attention allow a single model to generalize across heterogeneous workloads, eliminating the operational burden of per-application model maintenance. \ac{AgentEdge} established that simulation-validated multi-agent coordination achieves 78.3\% orchestration success rates (2.76$\times$ higher than single-agent baselines), enabling autonomous decision-making without trial-and-error exploration in production.

Together, these contributions transform orchestration from a reactive, human-supervised process into a proactive, self-managing capability. The research has produced five journal publications, three international conference papers in \ac{IEEE} venues, and one Elsevier book chapter, contributing validated methodologies and reproducible frameworks to the research community. As \ac{6G} networks evolve toward supporting thousands of distributed microservices under sub-millisecond latency constraints, the integrated pipeline of edge-deployable prediction, generalizable forecasting, and validated autonomous action provides a foundation for the zero-touch service orchestration that next-generation networks demand.

The transition from human-operated to autonomous service orchestration is no longer theoretical. This dissertation demonstrates it is achievable.
